\documentclass[11pt,a4paper]{article}
\usepackage[utf8]{inputenc}
\usepackage[T1]{fontenc}
\usepackage{geometry}
\usepackage{amsmath,amssymb,amsfonts}
\usepackage{booktabs}
\usepackage{array}
\usepackage{longtable}
\usepackage{hyperref}
\usepackage{enumitem}
\usepackage{xcolor}
\geometry{margin=1in}

\title{A Frozen-DINOv3 Query-Based Pipeline for Variable-Cardinality Palette Prediction}
\author{Project Report}
\date{\today}

\begin{document}
\maketitle

\begin{abstract}
This document describes the implemented supervised pipeline for palette prediction from JPG images. The system maps one image to an unordered set of RGB colors with variable cardinality using a frozen DINOv3 ViT-L/16 backbone, a DETR-style palette query decoder, and Hungarian bipartite matching. The report is written in a scientific-paper style and explains the method step by step, including data validation, preprocessing, feature extraction, mask handling, side-channel color features, loss design, evaluation metrics, inference, and training diagnostics.
\end{abstract}

\section{Problem Definition}

The task is to predict an unordered palette from a single JPG image:
\begin{equation}
I \mapsto \mathcal{P} = \{c_1, c_2, \dots, c_K\},
\end{equation}
where each color $c_i \in [0,255]^3$ is an RGB triplet and $K$ is image-dependent. Since palette colors have no semantic ordering, the target is a \emph{set}, not a sequence. Therefore, the learning objective must be permutation invariant with respect to the target row order in the annotation JSON.

The implemented architectural choice is:
\begin{center}
\texttt{JPG image} $\rightarrow$ \texttt{Frozen DINOv3 features} $\rightarrow$ \texttt{Query decoder} $\rightarrow$ \texttt{unordered RGB palette}.
\end{center}

\section{Data Sources and Pairing}

The dataset is formed by two aligned directories:
\begin{itemize}[leftmargin=2em]
    \item images: \texttt{E:/End-to-End/Data/jpg}
    \item palette annotations: \texttt{E:/End-to-End/Data/palettes}
\end{itemize}

Files are paired strictly by zero-padded stem:
\begin{equation}
\texttt{000000.jpg} \leftrightarrow \texttt{000000.json}.
\end{equation}

Each annotation is expected to contain:
\begin{itemize}[leftmargin=2em]
    \item \texttt{sample\_id}
    \item \texttt{num\_colors}
    \item \texttt{rgb}
    \item \texttt{hex}
\end{itemize}

\section{Dataset Validation and Filtering}

\subsection{Validation Rules}

The validation stage performs the following checks for every pair:
\begin{enumerate}[leftmargin=2em]
    \item Verify that the file stem is numeric.
    \item Verify that \texttt{sample\_id = int(stem)}.
    \item Verify that
    \begin{equation}
    \texttt{num\_colors} = |\texttt{rgb}| = |\texttt{hex}|.
    \end{equation}
    \item Verify that every RGB entry contains exactly three integers in $[0,255]$.
    \item Verify that every HEX string matches its RGB counterpart exactly.
    \item Detect duplicate RGB values inside the same palette.
\end{enumerate}

Originally, duplicate RGB values were only warnings. In the revised pipeline, duplicate colors are treated as invalid training targets because the target is mathematically a set and repeated elements create ambiguity.

\subsection{Training Cardinality Cap}

The training set is explicitly restricted to palettes with at most 12 colors:
\begin{equation}
K \leq 12.
\end{equation}

This policy was introduced after observing extreme outliers with very large palette cardinalities. Using a very large
\texttt{max\_num\_colors} would force the model to allocate many unnecessary negative queries and destabilize presence supervision.

Thus, the trainable subset is:
\begin{equation}
\mathcal{D}_{trainable} = \{(I, \mathcal{P}) \mid |\mathcal{P}| \leq 12\}.
\end{equation}

Samples above this threshold are not silently discarded; they are reported explicitly in the validation output.

\section{Preprocessing}

\subsection{Aspect-Preserving Resize and Letterboxing}

The image transform proceeds in four steps:
\begin{enumerate}[leftmargin=2em]
    \item Load the JPG as RGB.
    \item Resize the image while preserving aspect ratio so that the long side becomes 512.
    \item Place the resized content into a fixed $512 \times 512$ square canvas via letterboxing.
    \item Normalize the tensor using the exact mean and standard deviation required by \texttt{vitl16\_lvd}.
\end{enumerate}

Let the raw image be $I \in \mathbb{R}^{H \times W \times 3}$. After resize and letterbox, the model receives
\begin{equation}
\tilde{I} \in \mathbb{R}^{512 \times 512 \times 3}.
\end{equation}

\subsection{Valid-Pixel Mask}

The preprocessing stage also produces a valid-content mask
\begin{equation}
M_{pix} \in \{0,1\}^{512 \times 512},
\end{equation}
where
\begin{equation}
M_{pix}(x,y)=
\begin{cases}
1 & \text{if } (x,y) \text{ belongs to resized real image content},\\
0 & \text{if } (x,y) \text{ belongs to letterbox padding}.
\end{cases}
\end{equation}

This mask is critical because padded areas must not be interpreted as valid palette evidence.

\subsection{Raw Unnormalized Tensor}

Besides the normalized input tensor, the pipeline preserves the raw pre-normalization tensor. This is needed for the chromatic side-channel described later, because exact color statistics should be computed in raw image space rather than from standardized backbone input.

\section{Frozen DINOv3 Backbone}

The backbone is the official DINOv3 ViT-L/16 LVD model, referenced in the codebase as \texttt{vitl16\_lvd}. The backbone is frozen:
\begin{equation}
\frac{\partial \mathcal{L}}{\partial \theta_{DINO}} = 0.
\end{equation}

All training updates are restricted to the palette head.

\subsection{Multi-Level Intermediate Features}

The feature extractor obtains normalized intermediate patch tokens and the class token using the public DINOv3 API. The selected transformer block indices are
\begin{equation}
[5, 11, 17, 23].
\end{equation}

These should be understood as \emph{multi-level blocks}, not literally the last four blocks. The design rationale is that color-sensitive appearance cues may survive better when features are aggregated across depth instead of relying only on the final representation.

For each selected block $l$, the backbone returns:
\begin{equation}
T^{(l)} \in \mathbb{R}^{N \times C}, \qquad \text{CLS}^{(l)} \in \mathbb{R}^{C},
\end{equation}
where $N$ is the number of image patches and $C$ is the hidden dimension of the backbone.

The multi-level patch memory is formed by concatenation:
\begin{equation}
T_{fused} = \text{concat}\left(T^{(5)}, T^{(11)}, T^{(17)}, T^{(23)}\right) \in \mathbb{R}^{N \times 4C}.
\end{equation}

The final class token is taken from the deepest selected block:
\begin{equation}
\text{CLS}_{final} = \text{CLS}^{(23)}.
\end{equation}

\section{Patch-Space Masking}

Because the backbone operates at patch resolution, the pixel mask is converted into a patch mask:
\begin{equation}
M_{patch} \in \{0,1\}^{N}.
\end{equation}

Each patch is marked valid if most of its pixels correspond to true image content. This patch mask is then used as a transformer memory padding mask:
\begin{equation}
\texttt{memory\_key\_padding\_mask} = \neg M_{patch}.
\end{equation}

This ensures that letterbox padding does not influence:
\begin{itemize}[leftmargin=2em]
    \item cross-attention from queries to memory,
    \item pooled count features,
    \item auxiliary count evidence.
\end{itemize}

\section{Raw-Color Side-Channel}

DINOv3 is strong at structural and semantic representation learning, but exact RGB prediction may require more direct access to pixel-level chromatic statistics. Therefore, the implemented pipeline augments the DINO memory with a compact patch-wise raw-color descriptor.

For each valid patch, the side-channel computes:
\begin{itemize}[leftmargin=2em]
    \item mean RGB,
    \item RGB standard deviation,
    \item RGB minimum,
    \item RGB maximum,
    \item mean OKLab.
\end{itemize}

Let this descriptor be:
\begin{equation}
R_i \in \mathbb{R}^{d_r}
\end{equation}
for patch $i$. Then the final memory token for patch $i$ is formed as:
\begin{equation}
Z_i = W \left[ T_{fused,i} \, \Vert \, R_i \right],
\end{equation}
where $\Vert$ denotes concatenation and $W$ is a trainable projection.

This preserves the original backbone choice while allowing the head to ``read'' direct color evidence from the image.

\section{Query-Based Palette Decoder}

The head is DETR-like and uses a fixed number of learned queries:
\begin{equation}
Q = K_{max} + Q_{spare}.
\end{equation}

With the current training cap:
\begin{equation}
K_{max}=12, \qquad Q_{spare}=4, \qquad Q=16.
\end{equation}

Each query is decoded through a pre-normalized transformer decoder. Each decoded query emits:
\begin{itemize}[leftmargin=2em]
    \item RGB prediction $\hat{c}_q \in [0,1]^3$,
    \item presence logit $\hat{p}_q \in \mathbb{R}$.
\end{itemize}

The model additionally predicts the global palette count:
\begin{equation}
\hat{K} = \arg\max \texttt{count\_logits}.
\end{equation}

\section{Count Prediction}

Count prediction is not based only on CLS. Instead, the implemented count branch aggregates three signals:
\begin{equation}
f_{count} = \left[
\text{CLS}_{proj}
\; \Vert \;
\text{MeanPool}(Z, M_{patch})
\; \Vert \;
\text{AttnPool}(Z, M_{patch})
\right].
\end{equation}

This is motivated by the concern that rare colors may weakly affect CLS alone. Masked mean pooling and masked attention pooling provide additional evidence grounded in valid patch memory.

\section{Hungarian Matching}

Because the target is an unordered set, the model uses one-to-one Hungarian assignment between predicted queries and target colors.

For predicted query $q$ and target color $j$, the matching cost is:
\begin{equation}
\mathcal{C}_{qj}
=
\lambda_{rgb} \| \hat{c}_q - c_j \|_1
+
\lambda_{perc} d_{OKLab}(\hat{c}_q, c_j)
-
\lambda_{pres} \log \sigma(\hat{p}_q).
\end{equation}

The matching is detached from gradient flow and solved with Hungarian bipartite assignment.

\section{Loss Function}

After matching, the final training loss is composed of several terms.

\subsection{Matched RGB Regression}

For matched pairs $(q,j)$:
\begin{equation}
\mathcal{L}_{rgb}
=
\frac{1}{B}
\sum_{b=1}^{B}
\left(
\frac{1}{|\mathcal{M}_b|}
\sum_{(q,j)\in\mathcal{M}_b}
\text{SmoothL1}(\hat{c}_{b,q}, c_{b,j})
\right).
\end{equation}

The key design detail is that the loss is averaged \emph{inside each image first}, then across the batch. This prevents images with large palettes from dominating optimization.

\subsection{Perceptual Color Loss}

For the same matched pairs:
\begin{equation}
\mathcal{L}_{perc}
=
\frac{1}{B}
\sum_{b=1}^{B}
\left(
\frac{1}{|\mathcal{M}_b|}
\sum_{(q,j)\in\mathcal{M}_b}
d_{OKLab}(\hat{c}_{b,q}, c_{b,j})
\right).
\end{equation}

\subsection{Presence Loss}

Presence supervision is strongly imbalanced because many queries do not correspond to true palette entries. Therefore, the implementation uses a focal-style binary loss:
\begin{equation}
\mathcal{L}_{pres} = \text{FocalBCE}(\hat{p}, y_{pres}).
\end{equation}

Positive and negative presence losses are logged separately for diagnostics.

\subsection{Count Loss}

Count is supervised with categorical cross-entropy:
\begin{equation}
\mathcal{L}_{count}
=
\text{CE}(\texttt{count\_logits}, K).
\end{equation}

\subsection{Count Consistency Loss}

The model maintains a second implicit count estimate from the sum of presence probabilities:
\begin{equation}
\hat{K}_{pres} = \sum_q \sigma(\hat{p}_q).
\end{equation}

Consistency is encouraged with:
\begin{equation}
\mathcal{L}_{cons}
=
\left| \hat{K}_{pres} - K \right|.
\end{equation}

\subsection{Auxiliary Decoder Losses}

Intermediate decoder layers can also produce RGB and presence outputs. Auxiliary supervision is applied to those intermediate predictions, but count is not duplicated at every layer.

\subsection{Total Loss}

The final loss is:
\begin{equation}
\mathcal{L}
=
\lambda_{rgb}\mathcal{L}_{rgb}
+
\lambda_{perc}\mathcal{L}_{perc}
+
\lambda_{pres}\mathcal{L}_{pres}
+
\lambda_{count}\mathcal{L}_{count}
+
\lambda_{cons}\mathcal{L}_{cons}
+
\lambda_{aux}\mathcal{L}_{aux}.
\end{equation}

\section{Evaluation Metrics}

All evaluation is Hungarian-aligned and permutation invariant. Reported metrics include:
\begin{itemize}[leftmargin=2em]
    \item count accuracy,
    \item count MAE,
    \item matched RGB MAE,
    \item $\Delta E00$ mean, median, P90, P95,
    \item percentages with $\Delta E00 \le 1,2,5$,
    \item exact-cardinality rate,
    \item palette success rate,
    \item false-positive query rate,
    \item false-negative color rate.
\end{itemize}

\subsection{Count Agreement Metrics}

Two count pathways are compared explicitly:
\begin{equation}
\hat{K}_{head} = \arg\max \texttt{count\_logits},
\qquad
\hat{K}_{pres} = \text{round}\left(\sum_q \sigma(\hat{p}_q)\right).
\end{equation}

The pipeline reports:
\begin{itemize}[leftmargin=2em]
    \item count-head accuracy,
    \item presence-derived count accuracy,
    \item absolute disagreement $|\hat{K}_{head} - \hat{K}_{pres}|$.
\end{itemize}

\subsection{Unified Set Error}

To avoid selecting checkpoints based only on matched-color quality, a unified set-level error is used:
\begin{equation}
E_{set}
=
\frac{
\sum_i \min(\Delta E00_i, c)
+
c |K - \hat{K}|
}{
\max(K,\hat{K})
},
\end{equation}
with default penalty constant $c=10$.

This metric jointly reflects color quality and cardinality correctness and is used for:
\begin{itemize}[leftmargin=2em]
    \item best checkpoint selection,
    \item early stopping,
    \item model comparison.
\end{itemize}

\section{Training Loop}

The training loop uses:
\begin{itemize}[leftmargin=2em]
    \item AdamW,
    \item warmup plus cosine decay,
    \item automatic mixed precision,
    \item gradient clipping,
    \item optional gradient accumulation,
    \item deterministic seeding,
    \item checkpoint save/resume.
\end{itemize}

Extensive progress tracking was added through tqdm bars and epoch-level summaries.

\section{Early Stopping}

Early stopping was made deliberately more conservative by introducing:
\begin{itemize}[leftmargin=2em]
    \item larger patience,
    \item a minimum number of epochs before patience is consumed,
    \item a minimum improvement threshold $\delta$.
\end{itemize}

This reduces the chance of terminating training because of minor metric fluctuations.

\section{Inference}

At inference time, the model:
\begin{enumerate}[leftmargin=2em]
    \item predicts count logits,
    \item ranks queries by presence probability,
    \item keeps the highest-confidence unique RGB colors,
    \item rounds them deterministically to 8-bit RGB,
    \item serializes them with exact uppercase HEX values.
\end{enumerate}

A lightweight UI was added for checkpoint loading, image upload, palette-swatch visualization, and JSON inspection.

\section{Testing}

The current codebase includes focused tests for:
\begin{itemize}[leftmargin=2em]
    \item dataset pairing and validation,
    \item collation,
    \item head output shape,
    \item Hungarian matching,
    \item permutation invariance,
    \item presence supervision,
    \item count loss,
    \item patch-mask effect,
    \item unified set error,
    \item JSON formatting,
    \item checkpoint round-trip,
    \item preprocessing shape guarantees.
\end{itemize}

At the latest check, the suite status was:
\begin{equation}
17 \text{ passed}, \qquad 2 \text{ skipped}.
\end{equation}

\section{Current Status and Remaining Empirical Work}

The code path for the revised pipeline has been implemented, but a few empirical results still need to be regenerated after the latest changes:
\begin{itemize}[leftmargin=2em]
    \item full preflight report on the real dataset,
    \item actual tiny-overfit results on the revised architecture,
    \item refreshed quantitative report with current runtime metrics.
\end{itemize}

\section{Conclusion}

The implemented method now provides a stronger and more principled foundation for palette prediction than the initial version. The main improvements are:
\begin{itemize}[leftmargin=2em]
    \item trainable-set filtering with explicit reporting,
    \item mask-aware letterbox handling,
    \item raw-color side-channel fused with frozen DINO features,
    \item more stable matching and presence supervision,
    \item set-level checkpoint metric,
    \item stronger diagnostics and testing infrastructure.
\end{itemize}

In summary, the system remains faithful to the original high-level design --- frozen DINOv3 plus DETR-style set decoding --- while significantly improving robustness in data handling, optimization stability, and evaluation quality.

\end{document}
